\subsection{Simply connected spaces}

DEFINITION 3. --- A topological space is said to be simply connected if all its coverings are trivializable.

A simply connected space is connected. Indeed, if a space X is the disjoint union of two nonempty open sets U and V, the canonical injection of U into X is a covering which is not trivializable.

The empty space is simply connected. Every topological space reduced to a single point is simply connected.

Remarks. --- 1) Let B be a simply connected topological space and (E, p) a covering of B. If the space E is connected and nonempty, then p is a homeomorphism. Let, for example, G be a connected topological group and H a discrete subgroup of G, so that the canonical map p: G → G/H makes G a covering of G/H (I, p. 100, cor. 2 of th. 1). If the space G/H is simply connected, then p is a homeomorphism and H is the trivial subgroup.

\begin{enumerate}
\def\labelenumi{\arabic{enumi})}
\setcounter{enumi}{1}
\tightlist
\item
  Let B be a topological space in which every point has a simply connected neighborhood; then every covering of a covering of B is a covering of B. Indeed, consider a covering (E, p) of B and a covering (F, q) of E. Let us prove that (F, \(p \circ q\)) is a covering of B. Since the question is local in B, we may assume that the space B is simply connected and hence that E is a trivializable covering of B. Let V be a connected component of E. It is open and closed and \(p \mid V: V \to B\) is a homeomorphism (I, p. 69, proposition 1); consequently, the space V is simply connected. Every connected component W of \(q^{-1}(V)\) is open and closed in \(q^{-1}(V)\), hence in F, and the map q induces a homeomorphism from W onto V. The map \(p \circ q\) thus makes F a covering of B, trivializable (loc. cit.).
\end{enumerate}

PROPOSITION 3. --- Let B be a topological space. Let (E, p) be a covering of B and let y be a point of E. Let X be a simply connected topological space, let f: X → B be a continuous map and x a point of X such that f(x) = p(y). There exists a unique continuous lift g: X → E of the map f such that g(x) = y.

The maps \(g\) sought correspond bijectively (I, p. 9, prop. 3) to the continuous sections \(s\) of the covering \(\mathrm{pr}_1: \mathrm{X} \times_{\mathrm{B}} \mathrm{E} \to \mathrm{X}\) such that \(s(x) = (x, y)\). Such a section exists because the space X is simply connected; it is unique because the space X is connected (I, p. 34, cor. 1 of prop. 11).

Example 1. --- Let X be a simply connected topological space and let f be a continuous function from X into \(C^{*}\). Recall (I, p. 101, example 6) that the map \(z \mapsto e^{z}\) makes C a covering of \(C^{*}\). By prop. 3, there exists a continuous function \(h: X \to C\) such that \(f(x) = e^{h(x)}\) for all \(x \in X\). If \(h': X \to C\) is another continuous function such that \(f(x) = e^{h'(x)}\) for all \(x \in X\), there exists an integer \(q \in Z\) such that \(h' = h + 2\pi iq\).

Similarly, let \(n\) be an integer \(>0\); the map \(z \mapsto z^{n}\) makes \(C^{*}\) a covering of itself (I, p. 101, example 5). There therefore exists a continuous function \(k\) from X into \(C^{*}\) such that \(k(x)^{n} = f(x)\) for all \(x \in X\), for example the function \(k(x) = e^{h(x)/n}\). If \(k' : X \to C^{*}\) is another continuous function such that \(k'(x)^{n} = f(x)\) for all \(x \in X\), there exists an \(n\)th root of unity \(\mu \in C\) such that \(k' = \mu k\).

COROLLARY. — Let B be a topological space and let b be a point of B. Let E be a simply connected covering of B. For every point x of \(E_{b}\) the pointed space (E, x) is a universal covering of (B, b).

PROPOSITION 4. --- The product of two simply connected spaces, one of which is locally connected, is a simply connected space.

Let X and Y be simply connected spaces, and suppose that Y is locally connected. The space \(X \times Y\) is connected, since X and Y are (I, p. 124). Let \((Z, f)\) be a nonempty covering of \(X \times Y\); let us prove that it is trivializable. By Corollary 2 of I, p. 70, it suffices to prove that for every point \(z_{0}\) of Z, there exists a continuous section s of f such that \(s(f(z_{0})) = z_{0}\). Let therefore \(z_{0}\) be a point of Z. Set \((x_{0}, y_{0}) = f(z_{0})\). The subspace \(X \times \{y_{0}\}\) of X × Y is homeomorphic to X. Consequently, the covering deduced from Z over \(X \times \{y_{0}\}\) is trivializable and has a continuous section σ such that \(\sigma(x_{0}, y_{0}) = z_{0}\). Similarly, for every point x of X, there exists a continuous section \(\tau_{x}\) of f over \(\{x\} \times Y\) such that \(\tau_{x}(x, y_{0}) = \sigma(x, y_{0})\). For \((x, y) \in X \times Y\), set \(s(x, y) = \tau_{x}(x, y)\). The map s is a section of f, and one has \(s(x_{0}, y_{0}) = z_{0}\). Since the space Y is connected and locally connected, it follows from th. 1 of I, p. 35 that the map s is continuous.

PROPOSITION 5. --- A connected and locally connected topological space such that the intersection of any two connected open subsets is connected is a simply connected space.

Let B be such a topological space. Let \((\mathrm{E}, p)\) be a covering of B, let x be a point of E, and denote \(b = p(x)\). It is a matter of proving that there exists a continuous section s of p such that \(s(b) = x\) (cor. 2, I, p. 70). Let S be the set of pairs \((\mathrm{U}, s_{\mathrm{U}})\) where U is a connected open subset of B containing b and \(s_{U}\) is a continuous section of \(p_{\mathrm{U}}: p^{-1}(\mathrm{U}) \to \mathrm{U}\) such that \(s_{\mathrm{U}}(b) = x\). The set S is not empty (I, p. 34, prop. 10). Let \((\mathrm{U}, s_{\mathrm{U}})\) and \((\mathrm{V}, s_{\mathrm{V}})\) be elements of S. Then, \(s_{U} \mid U \cap V\) and \(s_{V} \mid U \cap V\) are continuous sections of \(p_{U \cap V}\) which take the value x at b. By hypothesis, \(U \cap V\) is connected; we therefore have \(s_{U} \mid U \cap V = s_{V} \mid U \cap V\) (I, p. 34, cor. 1 of prop. 11). Let A be the union of the open sets U as \((\mathrm{U}, s_{\mathrm{U}})\) runs through S, and let s: \(A \to E\) be the unique map such that \(s \mid U = s_{U}\) for every pair \((\mathrm{U}, s_{\mathrm{U}}) \in S\). The set A is open and connected (TG, I, p. 81, prop. 2), it contains b, and s is a continuous section of \(p_{A}\) such that \(s(b) = x\). It now suffices to prove that the set A is closed, which will imply that it is equal to B.

Let \(a\) be a point of B adherent to A, and let V be a connected open neighborhood of \(a\) such that the covering E is trivializable over V. There exists a point \(c\) in \(\mathrm{A} \cap \mathrm{V}\) and a continuous section \(s_{\mathrm{V}}\) of \(p_{\mathrm{V}}\) such that \(s_{\mathrm{V}}(c) = s(c)\). Let \(\mathrm{A}'\) be the open set \(\mathrm{A} \cup \mathrm{V}\). Since \(\mathrm{A} \cap \mathrm{V}\) is connected, there exists a continuous section \(s'\) of \(p_{\mathrm{A}'}\) which extends \(s\) and \(s_{\mathrm{V}}\) (I, p. 35, cor. 3 of prop. 11); the pair \((\mathrm{A}', s')\) belongs to \(\mathcal{S}\) and \(\mathrm{A}'\) is therefore contained in A. Consequently, \(a\) belongs to A, and A is closed.

COROLLARY. --- Every interval of the real line R is simply connected.

Indeed, the connected subspaces of R are the intervals (TG, IV, p. 8, th. 4) and the intersection of two intervals is an interval.

Example 2. --- The n-dimensional numerical space \(R^{n}\) (TG, VI, p. 1) is simply connected. Indeed, it is a product of simply connected and locally connected spaces (I, p. 126, prop. 4 and cor. of prop. 5 above). The same holds for every open or closed parallelepiped of \(R^{n}\). A parallelotope and an open or closed Euclidean ball in \(R^{n}\) are simply connected because they are homeomorphic to a block (TG, VI, p. 10, prop. 2 and I, p. 124, example).

PROPOSITION 6. --- Let X be a topological space. Let U₁ and U₂ be open (resp. closed) subspaces of X such that X = U₁ ∪ U₂. Suppose that U₁ and U₂ are simply connected and that their intersection U₁ ∩ U₂ is connected and nonempty. Then, the space X is simply connected.

Let (E, p) be a covering of X and let y be a point of E. It suffices to show that there exists a continuous section s of p such that \(s(p(y)) = y\) (I, p. 70, cor. 2 of prop. 1). Suppose for example that \(p(y)\) belongs to \(U_{1}\). There then exists a continuous section \(s_{1}: U_{1} \to E\) of \(p_{U_{1}}\) such that \(s_{1}(p(y)) = y\). Let x be a point of \(U_{1} \cap U_{2}\); there exists a continuous section \(s_{2}\) of \(p_{U_{2}}\) such that \(s_{2}(x) = s_{1}(x)\). By Corollary 3 of Proposition 11 of I, p. 34, there exists a continuous section s of p extending both \(s_{1}\) and \(s_{2}\).

Examples. --- 3) For \(n \geqslant 2\), the sphere \(\mathbf{S}_{n}\) (TG, VI, p. 10) is simply connected. Indeed, the sphere \(\mathbf{S}_{n}\) is the union of two closed hemispheres homeomorphic to \(\mathbf{B}_{n}\) whose intersection is homeomorphic to

\(S_{n-1}\) (cf. TG, VI, p. 12). For \(n \geqslant 2\), the sphere \(S_{n-1}\) is connected, whence the assertion.

On the other hand, the circle \(S_{1}\) is not simply connected. Indeed, the continuous map \(p: R \to S_{1}\) defined by \(p(\theta) = (\cos \theta, \sin \theta)\) makes R an infinite-degree covering of \(S_{1}\), which is connected and hence not trivializable.

\begin{enumerate}
\def\labelenumi{\arabic{enumi})}
\setcounter{enumi}{3}
\tightlist
\item
  Let E be a finite-dimensional vector space of dimension n over R, and let F be an affine subspace of E of codimension \(p \geqslant 3\). The set \(R^{p}\setminus\{0\}\) is homeomorphic to \(R \times S_{p-1}\) (TG, VI, p. 10, cor. 2), hence it is simply connected, since R and \(S_{p-1}\) are locally connected and simply connected (I, p. 126, prop. 4). The set \(E\setminus F\), homeomorphic to \(R^{n-p} \times (R^{p}\setminus\{0\})\), is therefore simply connected (loc. cit.).
\end{enumerate}

PROPOSITION 7. — Let X be a topological space. Let \(U_{1}\) and \(U_{2}\) be connected open (resp. closed) subspaces of X such that \(X = U_{1} \cup U_{2}\). If the space X is simply connected, then \(U_{1} \cap U_{2}\) is connected.

Let us set \(U = U_{1} \cap U_{2}\) and suppose for contradiction that the space U is not connected. We will construct a connected covering of X of infinite degree; such a covering is not trivializable.

By hypothesis, the set U is the union of two disjoint, nonempty sets A and B that are open (resp. closed) in X. For \(i \in \{1, 2\}\), let \(Y_i\) be the \(U_i\)-space \((\mathrm{U}_i \times \mathbf{Z}, \mathrm{pr}_1)\), and let \(Z_i\) be the subspace \(U \times \mathbf{Z}\) of \(Y_i\). The map \(h: Z_1 \to Z_2\) defined by

\[
h (x, n) = \left\{ \begin{array}{l l} (x, n) & \text {if } x\in A \text { and } n\in \mathbf {Z} \\ (x, n + 1) & \text {if } x\in B \text { and } n\in \mathbf {Z} \end{array} \right.
\]

is a homeomorphism. Let Y be the space obtained by gluing \(Y_{1}\) and \(Y_{2}\) along \(Z_{1}\) and \(Z_{2}\) by means of the homeomorphism h (TG, I, p. 17).

For \(i \in \{1,2\}\), the canonical map \(g_i\) of \(Y_i\) into Y is a homeomorphism of \(Y_i\) onto an open (resp. closed) subset of Y (loc. cit., prop. 9). For every integer \(n \in \mathbf{Z}\), the sets \(g_i(U_i \times \{n\})\), \(i \in \{1,2\}\), are connected; since A and B are nonempty, \(g_1(U_1 \times \{n\})\) meets \(g_2(U_2 \times \{n\})\) and \(g_2(U_2 \times \{n+1\})\). It follows that the space Y is connected (TG, I, p. 81, corollary).

For \(x \in \mathrm{U}\) and \(n \in \mathbf{Z}\), we have \((\mathrm{pr}_1 \circ h)(x, n) = x = \mathrm{pr}_1(x, n)\). There therefore exists a unique continuous map \(p \colon \mathrm{Y} \to \mathrm{X}\) such that \(p \circ g_i = \mathrm{pr}_1\) for \(i \in \{1, 2\}\). Let us prove that the X-space \((\mathrm{Y}, p)\) is a covering.

By construction, the fibers of the map p are homeomorphic to the discrete space Z.

For \(i \in \{1,2\}\), the map \(g_i\) defines, by passing to subspaces, an isomorphism of \(U_i\)-spaces from \(\mathrm{U}_i \times \mathbf{Z}\) onto \(p^{-1}(\mathrm{U}_i)\).

By definition of the space Y, there exists a unique map k from \((\mathrm{X}\setminus\mathrm{A})\times\mathbf{Z}\) into Y such that \(k(x,n)=g_{1}(x,n)\) for \(x\in(\mathrm{U}_{1}\setminus\mathrm{A})\times\mathbf{Z}\) and \(k(x,n)=g_{2}(x,n-1)\) for \(x\in(\mathrm{U}_{2}\setminus\mathrm{A})\times\mathbf{Z}\); it is an isomorphism of \((\mathrm{X}\setminus\mathrm{A})\)-spaces from \((\mathrm{X}\setminus\mathrm{A})\times\mathbf{Z}\) onto \(p^{-1}(\mathrm{X}\setminus\mathrm{A})\). Likewise, there exists a unique map \(k'\) from \((\mathrm{X}\setminus\mathrm{B})\times\mathbf{Z}\) into Y which coincides with \(g_{1}\) on \((\mathrm{U}_{1}\setminus\mathrm{B})\times\mathbf{Z}\) and with \(g_{2}\) on \((\mathrm{U}_{2}\setminus\mathrm{B})\times\mathbf{Z}\); it is an isomorphism of \((\mathrm{X}\setminus\mathrm{B})\)-spaces from \((\mathrm{X}\setminus\mathrm{B})\times\mathbf{Z}\) onto \(p^{-1}(\mathrm{X}\setminus\mathrm{B})\).

This proves that the X-space \((Y,p)\) is trivializable over the parts \(U_{1}\), \(U_{2}\), \(X\setminus A\), and \(X\setminus B\). We have \(U_{1} \cup U_{2} = X\); if \(U_{1}\) and \(U_{2}\) are open in X, this proves that \((Y,p)\) is a covering of X. The same holds when \(U_{1}\) and \(U_{2}\) are closed in X because then \(X\setminus A\) and \(X\setminus B\) are open subsets of X whose union is X. The proposition is thus proved.

COROLLARY. --- Let X be a simply connected topological space and let A be a connected subset of X. If the complement of A is connected, then its boundary is also connected.

Let \(X_{1}\) and \(X_{2}\) be the closures of A and of \(X\setminus A\), respectively. The sets \(X_{1}\) and \(X_{2}\) are closed and connected (TG, I, p. 81, prop. 1); we have \(X_{1} \cup X_{2} = X\) and their intersection \(X_{1} \cap X_{2}\) is equal to the boundary of A. It then suffices to apply Proposition 7.
% semantic-fragments: legacy-input-seam
\relax{}
